\chapter{Wstęp}
\label{chap:intro}
W~ostatnich latach sieci neuronowe
zyskują coraz bardziej na znaczeniu
dzięki dynamicznemu rozwojowi uczena maszynowego.
Są one obecnie używane jako jeden z~podstawowych
narzędzi służących do analizy danych.
Znajdują zastosowanie między innymi w~klasyfikacji,
regresji, analizie obrazów, znajdywaniu wzorców
czy przetwarzaniu języka naturalnego,
ponieważ sprawują się lepiej
niż klasyczne metody statystyczne.

Skuteczność modeli zależy jednak nie tylko
od dobrze dobranych danych uczących,
lecz również od optymalności topologii danej sieci
oraz odpowiednich parametrów uczenia.
Ich dobór jest procesem wymagającym doświadczenia,
jednakże nawet posiadanie odpowiednich
kompetencji nie gwarantuje w~praktyce
zaprojektowania zadowalającej struktury.
W~konsekwencji tego potrzeba licznych eksperymentów
z~różnymi modelami i~parametrami uczenia,
by zapewnić coraz lepiej działające sieci.

Określenie struktury sieci neuronowych
oraz hiperparametrów ma istotny wpływ
na szybkość modelu, zdolność do generalizacji
oraz na stabilność procesu uczenia.
Tworzą one jednak dużą przestrzeń możliwych rozwiązań.
Dodatkowo by upewnić się, że dana konfiguracja
jest odpowiednia należy przeprowadzić kilka razy
proces uczenia po nowej inicjalizacji wag,
ponieważ jest to proces losowy,
a~może znacząco wpływać na końcowe zdolności sieci.
Wskazuje to na potrzebę ograniczenia ręcznego
zmieniania parametrów i~testowania,
i~kieruje do automatyzacji tego procesu.

Dobieranie wyżej wymienionych własności można
rozpatrywać jako problem optymalizacyjny
o~wysokiej wymiarowości przestrzeni.
Dziedzina ta składa się zarówno ze zmiennych dyskretnych
(np. liczba neuronów, epok czy rodzaj funkcji aktywacji)
jak i~ciągłych (np. współczynnik uczenia).
Dodatkowo funkcja celu nie posiada postaci analitycznej,
a~jej wartość jest uzyskiwana dopiero
po przeprowadzeniu procesu uczenia,
który jest kosztowny obliczeniowo.

Klasyczne metody przeszukiwania
hiperparametrów obejmują przeszukiwanie
siatki (ang. \textit{grid search})
oraz losowe przeszukiwanie
(ang. \textit{random search}).
Pierwsze z~nich polega na sprawdzeniu wszystkich
możliwych kombinacji z~uprzednio zdefiniowanego zbioru,
co będąc proste implementacyjnie jest również
wykładniczo rosnące w~kosztach obliczeniowych
wraz z~wzrostem liczby parametrów.
Druga w~wielu przpadkach okazuje się bardziej efektywna,
natomiast wymaga przeprowadzania dużej liczby
procesów uczenia dla tej samej architektury.
W~nowszej literaturze pojawiają się
metody optymalizacji bayesowskiej,
których to implementacja jest bardziej złożona,
a~skuteczność jest wrażliwa na przyjęte założenia.

Kolejną trudością jest otymalizacja architektury
sieci neuronowej, której to parametry
są dyskretne i~mają zmienną wymiarowość.
Utrudnia to zastosowanie metod gradientowych
co skłania ku wykorzystaniu metod,
które nie wymagają ciągłości
ani różniczkowalności funcji celu. 

Taką metodą jest między innymi algorytm genetyczny,
należący do rodziny algorytmów ewolucyjnych.
Algorytm ten dobrze radzi sobie
z~dużymi przestrzeniami rozwiązań,
dzięki stosowaniu zainspirowanych
z~natury mechanizmów takich jak
krzyżowanie, mutacja i~selekcja.
Klasycznie każde rozwiązanie jest reprezentowane
w~postaci kodu genetycznego, który jest modyfikowny
przez operatory genetyczne.

Ich zaletą jest zdolność do równoległego
przeszukiwania wielu obszarów przestrzeni rozwiązań,
co zmniejsza ryzyko utknięcia w~lokalnym ekstremum.
Sprawdzają się również dla mieszanych reprezentacji
oraz w~sytuacjach gdy nie istnieje
lub nie jest znany gradient funkcji celu
ani jej analityczna postać.
To sprawia, że idealnie pasują
do charakterystyki problematyki optymalizacji
architektury sieci neuronowej.

\section*{Cel}
Celem niniejszej pracy jest opracowanie biblioteki
w~języku Python umożliwiającej automatyczną
optymalizację architektury wielowarstwowego
perceptronu wraz z~jej parametrami uczenia
za pomocą algorytmu genetycznego.
Biblioteka ma umożliwić użytkownikowi przekazywanie
danych w~postaci tablic pakietu \texttt{numpy},
a~następnie po przeprowadzeniu procesu ewolucyjnego
powinno zwracać pełny obraz procesu wraz
z~podstawową statystyką.

\section*{Cele szczegółowe}
Pierwszym krokiem jest opracowanie struktury
kodu genetycznego umożliwiającego generowanie z~nich,
w~spójny sposób, architektury sieci neuronowej
oraz parametrów uczenia.

Kolejnym kamieniem milowym jest zaprojektowanie systemu
w~sposób umożliwiający łatwe rozszerzanie zakresu
przestrzeni rozwiązań oraz rozwijanie biblioteki. 

Kolejnym jest zawarcie w~kodzie nauczania
sieci neuronowych, które będzie względnie niezależne
od wybranej przez użytkownika platformy
(tensorflow, pytorch).

Ostatnim krokiem jest pełna archiwizacja prób
do wybranego folderu wraz z~nauczonymi wagami
z~każdego uczenia, by umożliwić użytkownikowi
wybranie już gotowego modelu
na podstawie dostarczonych danych,
bez następnego uczenia.

\section*{Pytania badawcze}
W~pracy podjęto próbę odpowiedzi
na następujące pytania badawcze:
\begin{nospitemize}
  \item Czy zastosowanie algorytmu genetycznego
    pozwala uzyskać modele o~jakości wyższej
    niż modele wybrane przez stochastyczną
    metodę przeszukiwania przestrzeni?
  \item Jak parametry algorytmu genetycznego wpływają
    na skuteczność procesu optymalizacji?
  \item Jak zmienia się czas trwania optymalizacji
    w~zależności od rozmiaru populacji
    i~liczby pokoleń?
\end{nospitemize}

\section*{Zakres}
Program powininen optymalizować
architekturę wielowarstwowych perceptronów
wraz z~jej parametrami uczenia
dla danych podanych przez użytkownika.
Przez architekturę,
rozumie się liczbę warstw,
liczbę neuronów w~każdej z~warstw
oraz funkcje aktywacji dla każdej z~warstw.
Parametry uczenia obejmują natomiast
rodzaj optymalizatora,
współczynnik uczenia, liczbę epok
oraz wartość wielkości wsadu.

Dane podawane przez użytkownika
powinny być przyjmowane w~postaci
tabel pakietu \texttt{numpy}.
Dwie tabele obowiązkowe -- treningowa oraz testowa.
Jedna tabela opcjonalna -- walidacyjna.

Przez optymalizację rozumie się znajdowanie
architektury sieci oraz parametrów jej uczenia
dla konkretnych danych dostarczonych przez użytkownika,
za pomocą algorytmu genetycznego.
Implementacja algorytmu genetycznego obejmuje
reprezentację osobników, mechanizm selekcji,
operatory krzyżowania i~mutacji
oraz procedurę oceny przystosowania.

Każdy osobnik ma przechowywać
jednoznaczny kod genetyczny, to znaczy,
że wybrany genotyp w~sposób
deterministyczny ma doprowadzić
do danej architektury
i~parametrów uczenia.

Przeprowadzenie eksperymentów obejmuje porównanie
określonych parametrów dla algorytmu genetycznego,
by wybrać dla niego domyślne parametry oraz sprawdzenie
czy automatycznie dobrane sieci są lepszej jakości
jak stworzone przez metodę stochastyczną.
Parametry obejmują liczbę generacji i~osobników
w~populacji oraz współczynnik krzyżowania i~mutacji.

\section*{Ograniczenia}
W~pracy rozważane są wyłącznie
sieci \acrshort{pMLP}.
Nie uwzględnia się ogólnego ich typu \acrshort{MLP},
sieci konwolucyjnych, rekurencyjnych
ani modeli z~połączeniami rezydualnymi.
Powoduje to, że sieci stworzona
w~wyniku działania biblioteki
może być ona dalej optymalizowana.

Uczenie sieci neuronowej jest operacją zużywającą
dużą ilość mocy obliczeniowej.
Dodatkowo ponieważ proces uczenia sieci zależy
od pierwotnie wylosowanych wag,
to dla uzyskania większej pewności
co do uzyskanego wyniku przystosowania należy
proces uczenia wykonać kilkakrotnie.
Powoduje to, że działanie programu
jest znacząco uzależnione
od mocy obliczeniowej urządzenia
wykonującego operacje
a~proces optymalizacji
poprzez algorytm genetyczny
jest stosunkowo długi.

\section*{Teza}
W~niniejszej pracy przyjęto tezę,
że wykorzystanie algorytmu genetycznego
do optymalizacji architektury sieci neuronowej
i~jej parametrów uczenia umożliwia uzyskanie modeli
o~jakości porównywalnej lub lepszej
niż rozwiązania projektowane manualnie
przy jednoczesnym ograniczeniu
nakładu pracy użytkownika.

\section*{Struktura pracy}
Niniejsza praca składa się
z~\total{chapter} rozdziałów.
Obecny rozdział opisuje cele pracy,
pytania badawcze, zakres, ograniczenia,
tezę oraz w~niniejszej sekcji
znajduje się informacja o~strukturze dokumentu.

\Cref{chap:theory} zawiera teorię
algorytmów genetycznych oraz sieci neuronowych,
z~szczególnym uwzględnieniem \acrshort{MLP}.
Znajduje się w~nim również przegląd wybranych
poprzednich rozwiązań optymalizujących sieci
za pomocą algorytmów ewolucyjnych.

\Cref{chap:prob_and_req} zawiera analizę problematyki
podjętej w~niniejszym dokumencie i~wymagań dla biblioteki,
pod względem tego co biblioteka powinna udostępniać,
jak winna się odbywać komunikacja
i~co biblioteka powinna w~ogólnym zarysie wykonywać.

\Cref{chap:sol_and_impl} zawiera opis projektu rozwiązania
i~implementacji. Przedstawia, co biblioteka powinna
udostępniać i~w~jaki sposób działa implementacja.

\Cref{chap:test_and_result} opisuje metodykę eksperymentów,
a~następnie pokazuje jakie są ich wyniki i~wstępnie
je omawia, pod względem charakterystyk wykresów.
